\documentclass[11pt, a4paper]{article}

% ── Encoding & fonts ──────────────────────────────────────────────
\usepackage[utf8]{inputenc}
\usepackage[T1]{fontenc}
\usepackage{lmodern}

% ── Page layout ───────────────────────────────────────────────────
\usepackage[margin=2.5cm]{geometry}

% ── Math ──────────────────────────────────────────────────────────
\usepackage{amsmath}
\usepackage{amssymb}
\usepackage{amsthm}
\usepackage{mathtools}

% ── Theorem environments ──────────────────────────────────────────
\newtheorem{theorem}{Theorem}[section]
\newtheorem{lemma}[theorem]{Lemma}
\newtheorem{proposition}[theorem]{Proposition}
\newtheorem{corollary}[theorem]{Corollary}
\newtheorem{definition}[theorem]{Definition}
\newtheorem{remark}[theorem]{Remark}

% ── Graphics & figures ────────────────────────────────────────────
\usepackage{graphicx}
\usepackage{float}
\usepackage{caption}
\usepackage{subcaption}

% ── Tables ────────────────────────────────────────────────────────
\usepackage{booktabs}
\usepackage{array}

% ── Algorithms ────────────────────────────────────────────────────
\usepackage{algorithm}
\usepackage{algpseudocode}

% ── Colors ────────────────────────────────────────────────────────
\usepackage{xcolor}
\definecolor{linkblue}{RGB}{0, 70, 140}

% ── Hyperlinks ────────────────────────────────────────────────────
\usepackage[
    colorlinks=true,
    linkcolor=linkblue,
    citecolor=linkblue,
    urlcolor=linkblue
]{hyperref}

% ── Bibliography ──────────────────────────────────────────────────
\usepackage[
    backend=biber,
    style=alphabetic,
    sorting=ynt
]{biblatex}
\addbibresource{references.bib}

% ── Misc ──────────────────────────────────────────────────────────
\usepackage{microtype}       % better text justification
\usepackage{parskip}         % space between paragraphs instead of indent
\usepackage{enumitem}        % better lists
\usepackage{bm}              % bold math
\usepackage{bbm}             % blackboard bold numbers

% ── Custom commands ───────────────────────────────────────────────
\newcommand{\R}{\mathbb{R}}
\newcommand{\Z}{\mathbb{Z}}
\newcommand{\N}{\mathbb{N}}
\newcommand{\E}{\mathbb{E}}
\newcommand{\F}{\mathbb{F}}
\newcommand{\norm}[1]{\left\|#1\right\|}
\newcommand{\abs}[1]{\left|#1\right|}
\newcommand{\inner}[2]{\langle #1, #2 \rangle}

% for your specific notation
\newcommand{\Wint}{W_{\text{int}}}
\newcommand{\What}{\hat{W}}
\newcommand{\SW}{S_W}
\newcommand{\Sx}{S_x}
\newcommand{\Sxz}{S_x^{(0)}}
\newcommand{\DeltaW}{\Delta W}
\newcommand{\zn}{z^{(n)}}
\newcommand{\zni}{z^{(n)}_{int}}
\newcommand{\znhat}{\hat{z}^{(n)}}
\newcommand{\Wn}{W^{(n)}}
\newcommand{\Wnhat}{\hat{W}^{(n)}}
\newcommand{\xn}{x^{(n)}}
\newcommand{\xnhat}{\hat{x}^{(n)}}
\newcommand{\SWn}{S_W^{(n)}}
\newcommand{\Sxn}{S_x^{(n)}}
\newcommand{\Sn}{S^{(n)}}
\newcommand{\Sxnpo}{S_x^{(n+1)}}
\newcommand{\en}{e^{(n)}}
\newcommand{\enpo}{e^{(n+1)}}
\newcommand{\xnpo}{x^{(n+1)}}
\newcommand{\znpo}{z^{(n+1)}}
\newcommand{\xnpoh}{\hat{x}^{(n+1)}}

% ── Title information ─────────────────────────────────────────────
\title{
    \vspace{-1cm}
    % {\large Midterm Report} \\[0.5em]
    {\LARGE Analysys of Quantization Error} \\[0.5em]
    % {\large for zkML Applications}
}
% \author{
%     Alex \\
%     \small Institution \\
%     \small \href{mailto:your@email.com}{your@email.com}
% }
% \date{\today}

% ══════════════════════════════════════════════════════════════════
\begin{document}
% ══════════════════════════════════════════════════════════════════

\maketitle

% \begin{abstract}
%     Write your abstract here. Typically 150--250 words summarizing
%     the problem, approach, and preliminary results.
% \end{abstract}

\tableofcontents
\newpage

% ──────────────────────────────────────────────────────────────────
% \section{Introduction}
% % ──────────────────────────────────────────────────────────────────

% Introduce the problem, motivation, and contributions.

% % ──────────────────────────────────────────────────────────────────
% \section{Background}
% % ──────────────────────────────────────────────────────────────────

% \subsection{Zero-Knowledge Machine Learning}

% \subsection{Neural Network Quantization}

% \subsection{Equivalence Verification}

% ──────────────────────────────────────────────────────────────────
\section{Error Analysis for Quantization in Linear Layers}
% ──────────────────────────────────────────────────────────────────

\subsection{Linear Layers}

We simplify linear layers to operation:
$$ z = W \cdot x $$

To add the bias in the error analysis isn't hard.

\subsection{Models}

\begin{itemize}
\item Description of how the model behaves through an inference
\item We simulate models containing a sequence of $L$ linear layers
\item $x^{(0)}$ the original input and $z^{(L)}$ final output
\item Quantization bring new steps
\begin{itemize}
	\item Quantization of the input
	\item After matrix mult, the output is rescaled to fit the chosen quantization range
\end{itemize}
\end{itemize}

\paragraph{Base model}
$$
x^{(0)} \xrightarrow{\text{linear}} z^{(0)} \xrightarrow{\text{=}} x^{(1)} \rightarrow z^{(1)} \rightarrow ... \rightarrow z^{(L)}
$$


\paragraph{Quantized model}
$$
x^{(0)}
\xrightarrow{\text{quantization}} x^{(0)}_{int}
\xrightarrow{\text{linear}} z^{(0)}_{int}
\xrightarrow{\text{rescaling}} x^{(1)}_{int}
\rightarrow z^{(1)}_{int} \rightarrow ... \rightarrow z^{(L)}_{int}
\xrightarrow{\text{dequantization}} \hat{z}^{(L)}
$$

\paragraph{Virtual quantized model converted back to floats for comparison with base model}
$$
x^{(0)}
\xrightarrow{\text{quantization}} \hat{x}^{(0)}
\xrightarrow{\text{linear}} \hat{z}^{(0)}
\xrightarrow{\text{rescaling}} \hat{x}^{(1)} \rightarrow \hat{z}^{(1)} \rightarrow ...
\rightarrow \hat{z}^{(L)}
$$



\subsection{Error sources}

Quantization introduce approximation error when:
\begin{itemize}
    \item quantizing input,
    \item quantizing parameters
    \item rescaling output of any matrix multiplication to keep them within quantization range
    \item clipping when value overflows the designated range
    \item non-linear layers approx (see in next section)
\end{itemize}

\subsubsection{Parameter Quantization}
The basic quantization relationship for parameters $W$:

\paragraph{Definition}
\begin{itemize}
    \item $W$ is the parameter of base model
    \item $\Wint$ is the quantized parameter
    \item $\What$ is the quantized parameter converted back to floats for comparison with base model $W$
    \item $\SW$ is the scaling factor for $W$ and is computed once during the quantization process
\end{itemize}

$$    \Wint = round\!\left(\frac{W}{\SW}\right)$$

$$
    \What = \Wint \cdot \SW = round\!\left(\frac{W}{\SW}\right) \cdot \SW
$$

\paragraph{Bounds}

Special property between base and quantized parameters.

\begin{equation}
 \forall i,j \quad \DeltaW_{ij} = \norm{W_{ij} - \What_{ij}}_{\infty}
%= | W - round\!\left(\frac{W}{\SW}\right) \cdot \SW |
\leq \frac{\SW}{2}
\end{equation}

Proof in Annex. (TODO)


\subsubsection{Input Quantization}

\paragraph{Definition}
\begin{itemize}
    \item this is a special case of rescaling
    \item $x$ is the input of base model
    \item $x_{int}$ is the quantized input
    \item $\hat{x}$ is the quantized input converted back to floats for comparison with base model $x$
    \item $\Sxz$ is the scaling factor for $x$ and is \emph{computed once during a calibration phase}
\end{itemize}

$$    x_{int} = round\!\left(\frac{x}{\Sxz}\right)$$

$$
    \hat{x} = x_{int} \cdot \Sx = round\!\left(\frac{x}{\Sxz}\right) \cdot \Sxz
$$


\paragraph{Bounds}

\begin{equation}
 \Delta x = \norm{ x - \hat{x}}_{\infty}
%= | W - round\!\left(\frac{W}{\SW}\right) \cdot \SW |
\leq \frac{\Sxz}{2}
\end{equation}

\paragraph{Proof}
This is a property of the rounding operation.
$$|a - round(\frac{a}{b}) \cdot b| \leq \frac{b}{2}$$


\subsubsection{Rescaling}
\begin{itemize}
	\item Happens after linear layers $z^{(n)} = W^{(n)} \cdot \xn$
	\item output $z^{(n)}$ is at scale $\SWn \cdot \Sxn$ and must scaling must be reset
	\item then scale back to designated range with $\Sxnpo$
	\item $\Sxnpo$ is the scaling factor for this linear layer and is \emph{computed once during a calibration phase}
\end{itemize}

$$ \xnpo_{int} = round\!\left(\frac{\zn_{int} \cdot \SWn \cdot \Sxn }{\Sxnpo}\right) = round\!\left(\frac{\znhat}{\Sxnpo}\right)$$

$$
    \xnpoh = \xnpo_{int} \cdot \Sxnpo = round\!\left(\frac{\znhat}{\Sxnpo}\right) \cdot \Sxnpo
$$

\paragraph{Bounds}

\begin{equation}
 \Delta \xn = \norm{\xn - \xnhat}_{\infty}
\leq \frac{\Sxn}{2}
\end{equation}

\paragraph{Note}
\begin{itemize}
	\item Input quantization is a special case of rescaling
	\item that's why we name rescale factor for input as $\Sx^{(0)}$
\end{itemize}


\subsubsection{Clipping}

\begin{itemize}
	\item Happens if some value overflows from designated range
	\item this should not happen if scale factors are chosen correctly
	\item but on some extreme inputs not accounted for it may happen
	\item in this study we have assumed that clipping does not occur and omitted it
\end{itemize}

\subsection{Computing error for a linear layer}

\begin{itemize}
	\item we compute error $\en$ for layer $n$
	\item in the quantized model a linear layer consists of the linear multiplication and the rescaling
	\item the error is the difference between base and the quantized model reverted back to floats
	\item at layer $n$: $\Delta \xnpo = |\xnpo - \xnpoh|$
	\begin{itemize}
    	\item  base model:
    	$x^{(n)} \xrightarrow{\text{linear}} z^{(n)} \xrightarrow{\text{=}} x^{(n+1)}$
    	\item quantized model: $x^{(n)}_{int}
    	\xrightarrow{\text{linear}} z^{(n)}_{int}
    	\xrightarrow{\text{rescaling}} x^{(n+1)}_{int} $
    	\item quantized model reversed to floats: $\hat{x}^{(n)}
    	\xrightarrow{\text{linear}} \hat{z}^{(n)}
    	\xrightarrow{\text{rescaling}} \hat{x}^{(n+1)} $
	\end{itemize}
	\item we name the error bound: $\en$
\end{itemize}

\paragraph{Result}
\begin{equation}
\boxed{
\Delta \xnpo \leq
\enpo =
\underbrace{\frac{\Sxn}{2} \cdot \norm{\Wn}_{\infty}}_{\text{accumulated error}} +
\underbrace{\frac{\SWn}{2} \cdot \norm{\xn}_1}_{\text{weight quant}} +
\underbrace{\frac{\SWn \cdot \Sxn}{4} \cdot d_{in}^{(n)}}_{\text{cross term}} +
\underbrace{\frac{\Sxnpo}{2}}_{\text{rescaling}}
}
\end{equation}
$\text{with } d_{in}^{(n)} \text{ the input dimension of layer } n$

\paragraph{Proof}
In annex.

\paragraph{Remark}
\begin{itemize}
\item Bound on error $\Delta \xnpo$ depends on previous input $x^{(n)}$
\item Some implementations have global quantization $\Sxn = \SWn = S$
\item In this case, we can simplify our bounds to

$$
\Delta \xnpo \leq
\frac{S}{2} \left( \norm{\xn}_1 + \norm{\Wn}_{\infty} + 1 \right)
\ \frac{S^{2}}{4} \cdot d_{in}^{(n)}
$$
$\text{with } d_{in}^{(n)} \text{ the input dimension of layer } n$

\end{itemize}

\subsection{Total error propagation}

\begin{itemize}
    \item The error from previous layer is amplified by matrix multiplication
\end{itemize}
The total error at layer $L$ satisfies:
\begin{equation}
\boxed{
    \norm{\Delta x^{(L)}}_{\infty} \leq
    \sum_{k=1}^{L}
    \left(\prod_{\ell=k+1}^{L} \norm{W^{(\ell)}}_\infty\right)
    \cdot e^{(k)}
    }
\end{equation}

\paragraph{Proof}
By induction. In annex. (TODO)

% ──────────────────────────────────────────────────────────────────
\section{Augmenting models with non-linear layers}
% ──────────────────────────────────────────────────────────────────

$$
x^{(0)}
\xrightarrow{\text{quantization}} \hat{x}^{(0)}
\xrightarrow{\text{linear}} \hat{z}^{(0)}
\xrightarrow{\text{rescaling}} \hat{x}^{(1)} \rightarrow \hat{z}^{(1)} \rightarrow ...
\rightarrow \hat{z}^{(L)}
$$
%
\subsection{ReLU}
ReLU does not bring new approximation since it works fine on integers.
\begin{itemize}
	\item Case 1: Both neurons are active -> Effective error pass through unchanged
	\item Case 2: Both neurons are inactive -> Effective error is killed and is 0
	\item Case 3: Activation flip (one active and one inactive) -> Effective error is reduced
\end{itemize}

\subsection{Others}

\begin{itemize}
	\item GELU, SiLU, Softmax, LayerNorm
	\item not computable on integers and need to be approximated
	\item this adds a new source of errors
	\item recent zkml works uses lookup tables for speed and precision improvement
	\item see NanoZK who boasts really low error on their lookups
\end{itemize}


\section{Annexes}

\subsection{Proof of error bounds on a single layer}

$$
\boxed{
\Delta \xnpo \leq
\underbrace{\Wn \cdot \frac{\Sxn}{2}}_{\text{accumulated error}} +
\underbrace{\frac{\SWn}{2} \cdot \xn}_{\text{weight quant}} +
\underbrace{\frac{\SWn \cdot \Sxn}{4}}_{\text{cross term}} +
\underbrace{\frac{\Sxnpo}{2}}_{\text{rescaling}}
}
$$

\subsubsection{Rescaling error}
\paragraph{Result}
$$\Delta \xnpo = |\Delta \zn - \Delta \Sxnpo|$$

\paragraph{Proof}
\begin{align*}
\Delta \xnpo &= \norm{\xnpo - \xnpoh}_{\infty} \\
&= \norm{\zn - \xnpo_{int} \cdot \Sxnpo}_{\infty} \\
&= \norm{\zn - round\!\left(\frac{\znhat}{\Sxnpo}\right) \cdot \Sxnpo }_{\infty} \\
&= \norm{\zn - round\!\left(\frac{\znhat}{\Sxnpo}\right) \cdot \Sxnpo }_{\infty} \\
&\intertext{Error due to rounding is represented by  $\Delta \Sxnpo$}  \\
&= \norm{\zn - \left(\frac{\znhat}{\Sxnpo} \cdot \Sxnpo + \Delta \Sxnpo\right) }_{\infty} \\
&= \norm{\zn - \znhat - \Delta \Sxnpo}_{\infty} \\
&= \norm{\Delta \zn - \Delta \Sxnpo}_{\infty} \\
&\leq \norm{\Delta \zn}_{\infty} + \norm{\Delta \Sxnpo}_{\infty} \\
&\leq \norm{\Delta \zn}_{\infty} + \frac{\Sxnpo}{2}
\end{align*}

\subsubsection{Accumulated error $\Delta \zn$}

\begin{itemize}
    \item we define $\Wnhat = \Wn + \Delta \Wn$
	\item we define $\xnhat = \xn + \Delta \xn$
\end{itemize}

\begin{align*}
\Delta \zn &= \norm{\zn - \znhat}_{\infty} \\
&= \norm{ \Wn \cdot \xn - \Wnhat \cdot \xnhat }_{\infty} \\
&= \norm{\Wn \cdot \xn - (\Wn + \Delta \Wn) \cdot (\xn + \Delta \xn)}_{\infty} \\
&= \norm{\Delta \Wn \cdot \xn + \Wn \cdot \Delta \xn + \Delta \Wn \cdot \Delta \xn}_{\infty} \\
&\leq \norm{\Delta \Wn \cdot \xn}_{\infty} + \norm{\Wn \cdot \Delta \xn}_{\infty} + \norm{\Delta \Wn \cdot \Delta \xn}_{\infty} \\
&\leq \max_i(|(\Delta \Wn \cdot \xn)_i|) + \max_i(|(\Wn \cdot \Delta \xn)_i|) + \max_i(|(\Delta \Wn \cdot \Delta \xn)_i|) \\
&\leq \max_i(|\sum_{j} \Delta \Wn_{ij} \cdot \xn_j)|) + \max_i(|\sum_{j} \Wn_{ij} \cdot \Delta \xn_i|) + \max_i(|\sum_{j} \Delta \Wn_{ij} \cdot \Delta \xn_i|) \\
&\leq \max_i(\sum_{j} |\Delta \Wn_{ij}| \cdot |\xn_j|) + \max_i(\sum_{j} |\Wn_{ij}| \cdot |\Delta \xn_i|) + \max_i(\sum_{j} |\Delta \Wn_{ij}| \cdot |\Delta \xn_i|) \\
&\leq \max_i(\sum_{j} |\Delta \Wn_{ij}| \cdot |\xn_j|) + \max_i(\sum_{j} |\Wn_{ij}| \cdot |\Delta \xn_i|) + \max_i(\sum_{j} |\Delta \Wn_{ij}| \cdot |\Delta \xn_i|) \\
&\leq \frac{\SWn}{2} \max_i(\sum_{j} |\xn_j|) + \frac{\Sxn}{2}\max_i(\sum_{j} |\Wn_{ij}|) + \sum_{j} \frac{\Sxn}{2} \frac{\SWn}{2} \\
&\leq \frac{\SWn}{2} \cdot \norm{\xn}_1 +  \frac{\Sxn}{2} \cdot \norm{\Wn}_{\infty} + \frac{\SWn \cdot \Sxn}{4} \cdot d_{in}^{(n)}
\end{align*}
$
\text{with } d_{in}^{(n)} \text{ the input dimension of layer } n
$

\subsubsection{Final result}
$$
\Delta \xnpo \leq
\frac{\SWn}{2} \cdot \norm{\xn}_1 +  \frac{\Sxn}{2} \cdot \norm{\Wn}_{\infty} + \frac{\SWn \cdot \Sxn}{4} \cdot d_{in}^{(n)} + \frac{\Sxnpo}{2}
$$
$\text{with } d_{in}^{(n)} \text{ the input dimension of layer } n$
% ──────────────────────────────────────────────────────────────────

% ──────────────────────────────────────────────────────────────────
\printbibliography
% ──────────────────────────────────────────────────────────────────

\end{document}
